\subsection*{Forma Canônica Racional}

A questão diante de nós nos leva ao coração da teoria de formas canônicas, um dos cumes da Álgebra Linear. Especificamente, trataremos da \textbf{Forma Canônica Racional}, também conhecida como Forma Normal de Frobenius. Esta forma é de suma importância, pois, ao contrário da Forma de Jordan, ela existe para qualquer operador linear sobre um corpo arbitrário $\mathbb{F}$, sem a necessidade de o polinômio característico se fatorar completamente em fatores lineares sobre $\mathbb{F}$.

Nosso problema trata de uma matriz real cujos autovalores são complexos. É precisamente neste cenário que a Forma Racional revela seu poder.

\subsection*{Enunciado do Problema}

Seja $A$ uma matriz real cujo polinômio característico é $p_A(x) = (x^2+1)^2(x^2+x+1)^2$.  
Determine as possíveis formas racionais de $A$.

\subsection*{Demonstração e Resolução}

\subsubsection*{Passo 1: Fundamentos Teóricos da Forma Racional}

A Forma Canônica Racional de uma matriz $A$ é uma matriz bloco-diagonal, onde cada bloco é uma \textit{matriz companheira} de um polinômio.

\textbf{Definição (Matriz Companheira):} Para um polinômio mônico $q(x)=x^k+a_{k-1}x^{k-1}+\cdots+a_1x+a_0$, sua matriz companheira $C(q(x))$ é a matriz $k\times k$ dada por:

\[
C(q(x)) =
\begin{pmatrix}
0 & 0 & \cdots & 0 & -a_0 \\
1 & 0 & \cdots & 0 & -a_1 \\
0 & 1 & \cdots & 0 & -a_2 \\
\vdots & \vdots & \ddots & \vdots & \vdots \\
0 & 0 & \cdots & 1 & -a_{k-1}
\end{pmatrix}
\]

Uma propriedade crucial de $C(q(x))$ é que seu polinômio característico (e também o minimal) é exatamente $q(x)$.

\textbf{Teorema Fundamental:} Toda matriz $A$ é semelhante a uma única Forma Canônica Racional, que é a soma direta (matriz bloco-diagonal) das matrizes companheiras de seus \textbf{divisores elementares}.

\[
R = \operatorname{diag}(C(q_1(x)), C(q_2(x)), \dots, C(q_r(x)))
\]

Os divisores elementares $q_i(x)$ são potências de polinômios irredutíveis sobre o corpo em questão (no nosso caso, $\mathbb{R}$).

A relação entre os divisores elementares e os polinômios da matriz é a seguinte:
\begin{enumerate}
\item \textbf{Polinômio Característico:} $p_A(x) = q_1(x)\cdot q_2(x)\cdots q_r(x)$.
\item \textbf{Polinômio Minimal:} $m_A(x)$ é o mínimo múltiplo comum (MMC) de todos os divisores elementares $\{q_1(x), \dots, q_r(x)\}$.
\end{enumerate}

\subsubsection*{Passo 2: Análise do Polinômio Característico}

O polinômio característico dado é $p_A(x) = (x^2+1)^2(x^2+x+1)^2$.

Primeiro, devemos encontrar seus fatores irredutíveis sobre $\mathbb{R}$.

\begin{itemize}
\item Seja $f_1(x)=x^2+1$. O discriminante é $\Delta=0^2-4(1)(1)=-4<0$. Logo, $f_1(x)$ é irredutível sobre $\mathbb{R}$.
\item Seja $f_2(x)=x^2+x+1$. O discriminante é $\Delta=1^2-4(1)(1)=-3<0$. Logo, $f_2(x)$ também é irredutível sobre $\mathbb{R}$.
\end{itemize}

Portanto, os divisores elementares de $A$ devem ser potências de $f_1(x)$ e $f_2(x)$.

\subsubsection*{Passo 3: Determinando os Conjuntos de Divisores Elementares}

As potências dos divisores elementares da forma $(x^2+1)^k$ devem somar $2$, e as potências dos fatores de $f_2(x)$ também devem somar $2$. Logo, há duas possibilidades para cada fator (partições do inteiro 2):

\textbf{Para $(x^2+1)^2$:}
\begin{itemize}
\item Partição 1: $2 \implies \{(x^2+1)^2\}$.
\item Partição 2: $1+1 \implies \{x^2+1, x^2+1\}$.
\end{itemize}

\textbf{Para $(x^2+x+1)^2$:}
\begin{itemize}
\item Partição A: $2 \implies \{(x^2+x+1)^2\}$.
\item Partição B: $1+1 \implies \{x^2+x+1, x^2+x+1\}$.
\end{itemize}

Combinando, temos $2\times 2 = 4$ possibilidades.

\subsubsection*{Passo 4: Construção das Possíveis Formas Racionais}

\textbf{Caso 1: Partições (1) e (A)}
\begin{itemize}
\item Divisores Elementares: $\{(x^2+1)^2, (x^2+x+1)^2\}$.
\item Polinômio Minimal: $m_A(x)=\operatorname{MMC}((x^2+1)^2,(x^2+x+1)^2)=(x^2+1)^2(x^2+x+1)^2=p_A(x)$.
\item Polinômios: $q_1(x)=(x^2+1)^2=x^4+2x^2+1$, $q_2(x)=(x^2+x+1)^2=x^4+2x^3+3x^2+2x+1$.
\item Forma Racional:
\[
R_1=\operatorname{diag}(C(q_1(x)),C(q_2(x)))
\]
\[
R_1=\begin{pmatrix}
0&0&0&-1&0&0&0&0\\
1&0&0&0&0&0&0&0\\
0&1&0&-2&0&0&0&0\\
0&0&1&0&0&0&0&0\\
0&0&0&0&0&0&0&-1\\
0&0&0&0&1&0&0&-2\\
0&0&0&0&0&1&0&-3\\
0&0&0&0&0&0&1&-2
\end{pmatrix}
\]
\end{itemize}

\textbf{Caso 2: Partições (1) e (B)}
\begin{itemize}
\item Divisores Elementares: $\{(x^2+1)^2, x^2+x+1, x^2+x+1\}$.
\item Polinômio Minimal: $m_A(x)=\operatorname{MMC}((x^2+1)^2, x^2+x+1)=(x^2+1)^2(x^2+x+1)$.
\item Forma Racional:
\[
R_2=\operatorname{diag}(C((x^2+1)^2),C(x^2+x+1),C(x^2+x+1))
\]
\[
R_2=\begin{pmatrix}
0&0&0&-1&0&0\\
1&0&0&0&0&0\\
0&1&0&-2&0&0\\
0&0&1&0&0&0\\
0&0&0&0&0&-1\\
0&0&0&0&1&-1
\end{pmatrix}
\]
\end{itemize}

\textbf{Caso 3: Partições (2) e (A)}
\begin{itemize}
\item Divisores Elementares: $\{x^2+1,x^2+1,(x^2+x+1)^2\}$.
\item Polinômio Minimal: $m_A(x)=\operatorname{MMC}(x^2+1,(x^2+x+1)^2)=(x^2+1)(x^2+x+1)^2$.
\item Forma Racional:
\[
R_3=\operatorname{diag}(C(x^2+1),C(x^2+1),C((x^2+x+1)^2))
\]
\[
R_3=\begin{pmatrix}
0&-1&0&0&0&0&0&0\\
1&0&0&0&0&0&0&0\\
0&0&0&-1&0&0&0&0\\
0&0&1&0&0&0&0&0\\
0&0&0&0&0&0&0&-1\\
0&0&0&0&1&0&0&-2\\
0&0&0&0&0&1&0&-3\\
0&0&0&0&0&0&1&-2
\end{pmatrix}
\]
\end{itemize}

\textbf{Caso 4: Partições (2) e (B)}
\begin{itemize}
\item Divisores Elementares: $\{x^2+1,x^2+1,x^2+x+1,x^2+x+1\}$.
\item Polinômio Minimal: $m_A(x)=\operatorname{MMC}(x^2+1,x^2+x+1)=(x^2+1)(x^2+x+1)$.
\item Forma Racional:
\[
R_4=\operatorname{diag}(C(x^2+1),C(x^2+1),C(x^2+x+1),C(x^2+x+1))
\]
\[
R_4=\begin{pmatrix}
0&-1&0&0&0&0&0&0\\
1&0&0&0&0&0&0&0\\
0&0&0&-1&0&0&0&0\\
0&0&1&0&0&0&0&0\\
0&0&0&0&0&-1&0&0\\
0&0&0&0&1&-1&0&0\\
0&0&0&0&0&0&0&-1\\
0&0&0&0&0&0&1&-1
\end{pmatrix}
\]
\end{itemize}

\subsection*{Conclusão}

Existem exatamente \textbf{quatro} classes de semelhança para uma matriz real $A$ com o polinômio característico dado. Cada classe é representada por uma das quatro Formas Canônicas Racionais $R_1,R_2,R_3,R_4$.  
A forma específica de $A$ dependerá de seu polinômio minimal, que pode ser um de quatro polinômios distintos.